\chapter{Ontwikkeling}%
\label{ch:ontwikkeling}

In dit hoofdstuk wordt de technische realisatie van de applicatie beschreven, waarbij de conceptuele ontwerpen uit het voorgaande hoofdstuk worden vertaald naar een functionele Proof of Concept voor de Meta Quest 3. De ontwikkeling is uitgevoerd binnen de kaders van Design Science Research, waarbij de focus ligt op de creatie van een technologisch robuust instrument dat didactische waarde combineert met technische performantie \cite{Hevner2004, Peffers2007}.

\section{Technologische stack en ontwikkelomgeving}

De XR SkillSim-applicatie is technisch gerealiseerd binnen de Unity-engine (versie 2022.3 LTS). De keuze voor een Long Term Support-versie is essentieel voor de stabiliteit van de Proof of Concept en garandeert compatibiliteit met de gebruikte Software Development Kits (SDK's) gedurende de gehele ontwikkelcyclus.

\subsection{Engine en Platform}

Het project is gebouwd voor de Meta Quest 3, een standalone XR-headset die draait op een Android-gebaseerd besturingssysteem. De ontwikkeling is uitgevoerd binnen de OpenXR-standaard. Dit biedt een abstractielaag bovenop de hardware-specifieke API's, waardoor de applicatie gebruikmaakt van Action Maps voor input in plaats van directe hardware-mapping. Dit is een strategische keuze met het oog op de schaalbaarheid: hierdoor kan de software met minimale aanpassingen worden geport naar andere apparaten binnen de 'HOGENT XR Academy'.

Voor de interactielaag is de Meta XR All-in-One SDK geïntegreerd in combinatie met de Unity XR Interaction Toolkit. Deze stack biedt de nodige componenten voor zowel controller-gebaseerde interactie als hand-tracking, wat cruciaal is voor de medische accuraatheid van de simulatie.

\subsection{Render Pipeline en Optimalisatie}

Vanwege de mobiele chipset (Snapdragon XR2 Gen 2) van de Meta Quest 3 is de Universal Render Pipeline ingezet. In tegenstelling tot de Built-in Render Pipeline maakt URP gebruik van een Single-pass Forward Renderer, wat de overhead op de CPU en GPU aanzienlijk vermindert.

De volgende technische optimalisaties zijn toegepast om de performantie te waarborgen:
\begin{itemize}
    \item Single Pass Instanced Rendering: Hierbij wordt de volledige scène in één keer naar een double-wide render target geschreven, waardoor het aantal draw calls bijna wordt gehalveerd vergeleken met multi-pass rendering.
    \item Foveated Rendering: Door gebruik te maken van de OVR Manager is Fixed Foveated Rendering geactiveerd. Dit verlaagt de resolutie van de pixels in de periferie van het gezichtsveld, waardoor rekenkracht vrijkomt voor de kern van de simulatie.
    \item Batching: Statische objecten in de klinische omgeving zijn geconfigureerd voor Static Batching, terwijl voor dynamische objecten (zoals de AED-onderdelen) GPU Instancing wordt toegepast om de overhead van de render-pipeline te minimaliseren.
\end{itemize}

\begin{table}[htbp]
    \centering

    \label{tab:tech_stack}
    \begin{tabular}{|l|l|}
        \hline
        \textbf{Component} & \textbf{Technologie} \\ \hline
        Game Engine & Unity 6 \\ \hline
        XR Standaard & OpenXR \\ \hline
        Render Pipeline & Universal Render Pipeline (URP) \\ \hline
        Scripting Language & C\# 9.0 (.NET Standard 2.1) \\ \hline
        SDK's & Meta XR Core SDK, Interaction SDK, XR Interaction Toolkit \\ \hline
        Serialization & Newtonsoft.Json (Json.NET) \\ \hline
    \end{tabular}
    \caption{Overzicht van de technologische stack.}
\end{table}

\section{Architecturale implementatie van de simulatie-logica}

De kern van XR SkillSim wordt gevormd door een robuuste logische laag die verantwoordelijk is voor het afdwingen van de medische protocollen. Om te vermijden dat de broncode vervalt in een onbeheersbare if-else-structuur, is gekozen voor een implementatie gebaseerd op het State Design Pattern.

\subsection{State Pattern in C\#: De ScenarioManager}

De ScenarioManager fungeert als een Singleton die de globale status van de simulatie beheert. De feitelijke logica van elke stap in het reanimatieproces is echter gedelegeerd naar specifieke State-objecten die de IScenarioState interface implementeren.

\begin{listing}[ht]
    \begin{minted}[frame=lines, framesep=2mm, baselinestretch=1.2, fontsize=\small]{csharp}
        public interface IScenarioState 
        {
            void EnterState(ScenarioManager manager);
            void UpdateState(ScenarioManager manager);
            void ExitState(ScenarioManager manager);
            void HandleInteraction(InteractionEvent e);
        }
    \end{minted}
    \caption{De IScenarioState interface en de basisstructuur van de ScenarioManager.}
    \label{listing:scenario_state}
\end{listing}

Door deze interface te gebruiken, hoeft de ScenarioManager niet te weten welke specifieke fase van de reanimatie actief is. Hij roept simpelweg de methoden van de huidige status aan. Dit maakt het systeem uiterst flexibel. Het toevoegen van een nieuwe medische handeling vereist enkel het aanmaken van een nieuwe klasse die deze interface implementeert.

\subsection{Event-Driven communicatie en ontkoppeling}

Om de verschillende subsystemen zoals de UI, de audio-instructies en de datalogging gescheiden te houden van de scenario-logica, is een Event-Driven architectuur opgezet met behulp van C\# Actions. 

In plaats van dat de ScenarioManager direct verwijst naar de UIManager, vuurt hij globale events af zoals OnStateChanged of OnActionFailed. Subsystemen abonneren zich op deze events. Dit verhoogt de modulariteit aanzienlijk. De LoggingManager kan bijvoorbeeld alle events opvangen en opslaan zonder dat de ScenarioManager zich bewust hoeft te zijn van het bestaan van een log-systeem.

\subsection{De State Machine flow en transities}

De overgang tussen verschillende fasen van de reanimatie wordt technisch afgehandeld via een TransitionTo-methode in de ScenarioManager. 

Wanneer aan de voorwaarden van een status is voldaan, bijvoorbeeld wanneer de student succesvol de hulpdiensten heeft gebeld, roept de actieve state deze methode aan. De TransitionTo-logica voert de volgende stappen asynchroon uit:
\begin{enumerate}
    \item De ExitState() methode van de huidige status wordt aangeroepen voor het opruimen van tijdelijke triggers of timers.
    \item De nieuwe status wordt geïnstantieerd en toegewezen als de actieve status.
    \item De EnterState() methode van de nieuwe status wordt uitgevoerd, waarbij events worden afgevuurd die de UI en audio-instructies voor de volgende stap initialiseren.
\end{enumerate}

Deze aanpak garandeert een deterministische flow van het scenario, waarbij de applicatie altijd in een bekende en gevalideerde status verkeert.

\section{Ontwikkeling van het interactiesysteem}

De interactielaag vormt de brug tussen de hardware-sensoren van de Meta Quest 3 en de logische State Machine van de simulatie. In plaats van standaard componenten te gebruiken, zijn er specifieke interactoren ontwikkeld die voldoen aan de eisen van medische precisie.

\subsection{Custom Interactors: Velocity-based Shake Interactor}

Voor de bewustzijnscontrole is een eenvoudige aanraking onvoldoende om een medische handeling te valideren. Er is een custom script ontwikkeld dat de lineaire snelheid en versnelling van de controllers of hand-meshes monitort. De interactor definieert een volume rond de schouders van de patiënt. Zodra de handen van de gebruiker dit volume betreden, analyseert het script de bewegingsvectoren in real-time.

De handeling wordt pas als succesvol gemarkeerd wanneer er binnen een tijdvenster van twee seconden een opeenvolging van tegengestelde vectoren wordt gedetecteerd die een drempelwaarde van 1.5 m/s overschrijden. Dit simuleert het krachtig schudden van de schouders.

\begin{listing}[ht]
    \begin{minted}[frame=lines, framesep=2mm, baselinestretch=1.2, fontsize=\small]{csharp}
    // Controleer of de beweging krachtig genoeg is
    if (handVelocity.magnitude > velocityThreshold) {
        if (Vector3.Dot(handVelocity.normalized, lastVelocityDir) < -0.5f) {
            shakeCount++;
            lastVelocityDir = handVelocity.normalized;
        }
    }
    \end{minted}
\caption{Logica voor de detectie van de schudbeweging.}
\label{listing:interaction_zones}
\end{listing}

\subsection{Spatial Logic en Snap Zones: Sockets en Tags}

Om medische objecten, zoals de AED-elektroden, te valideren, wordt gebruikgemaakt van een systeem van Sockets en Tags. Elke socket op de patiënt-mesh heeft een specifieke identifier die correspondeert met de anatomische locatie.

De validatielogica is gebaseerd op drie parameters:
\begin{itemize}
    \item Tag-filtering: Alleen objecten met de juiste medische tag kunnen interageren met de socket.
    \item Distance threshold: De elektrode snapt pas naar de torso wanneer deze binnen een straal van 5 cm van de socket komt.
    \item Alignment check: Het systeem vergelijkt de rotatie-vector van de elektrode met de normaalvector van de patiënt-mesh om te garanderen dat de plakzijde correct naar het lichaam gericht is.
\end{itemize}

\subsection{Hand-tracking integratie: OVR Skeleton Mapping}

Een grote uitdaging in de ontwikkeling was het mappen van de Oculus Virtual Reality Skeleton data naar de interactielaag. Hand-tracking levert een complexe hiërarchie van joints. Om dit bruikbaar te maken voor medische interacties, is een abstractielaag gebouwd die de ruwe data vertaalt naar stabiele interactiepunten.

Er wordt een onzichtbare Interaction Hand gegenereerd die de posities van de Oculus Virtual Reality Skeleton volgt, maar de colliders vereenvoudigt tot geometrische vormen zoals capsules en spheres]. Voor handelingen zoals de ademhalingscontrole wordt gebruikgemaakt van Pose Detection. Hierbij controleert de software of de specifieke hoeken tussen de vingersegmenten overeenkomen met de vereiste handhouding om de kin omhoog te liften.

\begin{figure}[htbp]
    \centering
    \includegraphics[width=0.8\textwidth]{Colliders.jpg}
    \caption{Visualisatie van de technische interactiezones (Gizmos) en colliders op de patiënt-mesh in Unity.}
    \label{fig:interaction_zones}
\end{figure}

\section{Realisatie van het dataloggingsysteem}

Een cruciaal onderdeel van de ontwikkeling is het genereren van objectieve evaluatiegegevens, een behoefte die sterk naar voren komt in de literatuur over medische XR-training \autocite{Radianti2020VRHE, Kyaw2019VR}. Er is een gespecialiseerde LoggingManager ontwikkeld die volledig ontkoppeld is van de visuele representatie en zich uitsluitend richt op de data-integriteit.

\subsection{xAPI-concept en Data-architectuur}

Het loggingsysteem is architecturaal gebaseerd op de Experience API standaard. Dit stelt de applicatie in staat om elke gebruikersinteractie om te zetten in een gestructureerd Statement bestaande uit een Actor, Verb en Object. 

In de praktijk vertaalt de LoggingManager binnenkomende events van de ScenarioManager naar deze statements. Elke interactie, beslissing en tijdsafhankelijke variabele wordt vastgelegd met een nauwkeurigheid tot op de milliseconde. Wanneer een student bijvoorbeeld de AED activeert, wordt dit gelogd als: Student\_01 (Actor) Pressed (Verb) AED\_PowerButton (Object). Deze gestructureerde aanpak is essentieel voor de latere validatie van de prestaties tegenover de medische protocollen.

\subsection{JSON-serialisatie en Persistentie}

Om de framerate op de Meta Quest 3 stabiel te houden en de CPU-belasting tijdens kritieke simulatiemomenten te minimaliseren, worden de verzamelde gegevens gedurende de sessie in een List<T> gebufferd in het werkgeheugen. Pas bij het beëindigen van het scenario wordt deze lijst geserialiseerd naar een JSON-bestandsstructuur met behulp van de Newtonsoft.Json bibliotheek.

\begin{listing}[ht]
    \begin{minted}[frame=lines, framesep=2mm, baselinestretch=1.2, fontsize=\small]{csharp}
    {
        "timestamp": "2026-05-04T22:45:12.453Z",
        "actor": "Student_01",
        "verb": "Placed",
        "object": "AED_Pad_Apex",
        "result": {
            "success": true,
            "responseTime": 4.2
        },
        "context": "AEDPreparationState"
    }
    \end{minted}
\caption{Voorbeeld geserialiseerd log-fragment.}
\label{listing:logging}
\end{listing}

Het wegschrijven van deze data gebeurt asynchroon naar het lokale bestandssysteem van de headset via de Application.persistentDataPath. Door deze datagedreven benadering is het mogelijk om patronen in fouten te herkennen en belangrijke prestatie-indicatoren, zoals de hands-on time of de pauzes tussen acties, op een objectieve manier te onderzoeken \autocite{Hung2024VR, Padilha2019VR}.

\section{Optimalisatie en technische uitdagingen}

Gezien de hardwarematige beperkingen van een standalone headset zoals de Meta Quest 3, vormde de optimalisatie van de applicatie een substantieel onderdeel van het ontwikkelproces. Het doel was om een constante verversingssnelheid van 72 Hz te handhaven, aangezien fluctuaties in de framerate direct kunnen leiden tot simulator sickness \autocite{Kennedy1993SSQ}.

\subsection{Profiling en Debugging}

Om prestatie-bottlenecks te identificeren, is er gebruikgemaakt van de Unity Profiler in combinatie met de Oculus Performance Metrics Tool. Tijdens de eerste testfases vertoonden de physics-berekeningen van de hand-colliders pieken in de CPU-belasting. Door middel van de profiler werd vastgesteld dat de FixedUpdate-cyclus overbelast raakte door een te hoog aantal actieve Mesh Colliders op de patiënt-torso.

Deze bottleneck is technisch opgelost door:
\begin{itemize}
    \item De complexe mesh-colliders te vervangen door een samengesteld systeem van Primitive Colliders, wat de rekenlast voor de Physics Engine aanzienlijk verminderde.
    \item Scripts die niet essentieel zijn voor de actieve status van de State Machine tijdelijk te deactiveren via de ScenarioManager, waardoor onnodige Update-aanroepen werden geëlimineerd.
\end{itemize}

\subsection{Beperkingen van Standalone XR: GPU-optimalisatie}

De mobiele GPU van de Quest 3 stelt strikte limieten aan de polycount en het aantal draw calls. Om de visuele presentatie te optimaliseren zonder de stabiliteit in gevaar te brengen, zijn de volgende strategieën toegepast:

\begin{itemize}
    \item Draw Call Reduction: Door gebruik te maken van Texture Atlasing zijn meerdere texturen van medische objecten samengevoegd in één enkele material, waardoor de GPU minder state changes hoeft uit te voeren.
    \item LOD-systemen: Voor de medische omgeving is Level of Detail geïmplementeerd, waarbij de complexiteit van de meshes dynamisch wordt verlaagd op basis van de afstand tot de camera van de headset.
    \item Shader Optimization: Er is uitsluitend gebruikgemaakt van de URP Lit en Simple Lit shaders, die specifiek zijn geoptimaliseerd voor mobiele architecturen en geen complexe real-time reflections vereisen.
\end{itemize}

Ondanks de technische uitdagingen met betrekking tot de afwezigheid van fysieke tegendruk en de beperkte rekenkracht, resulteerden deze optimalisaties in een stabiele leeromgeving die de procedurele nauwkeurigheid van de reanimatietraining waarborgt.



























\section{Technische aansturing van de simulatie-logica}
De essentie van de applicatie wordt beheerd door een centrale State Machine die de voortgang van de Basic Life Support (BLS)-procedure monitort. Elke fase in het algoritme van de European Resuscitation Council wordt technisch beschreven als een specifieke status met de bijbehorende transitievoorwaarden \cite{ERC2025BLS}. Als een gebruiker een actie uitvoert, zoals het controleren van zijn bewustzijn, wordt dit door het interactiesysteem omgezet in een gebeurtenis die de State Machine voedt. Als de actie correct en binnen de vastgestelde tijd wordt uitgevoerd, zal er een overgang plaatsvinden naar de volgende fase van de simulatie.

De interactie-logica maakt gebruik van 'interactors' en 'interactables' uit de Unity XR Interaction Toolkit \cite{UnityDocs}. Voor medische procedures die nauwkeurigheid vereisen, zoals het juiste plaatsen van AED-elektroden, zijn ruimtelijke triggers en collision detection toegepast. De applicatie monitort de coördinaten van de virtuele handen in relatie tot de anatomische referentiepunten op de torso van de patiënt. Deze technische validatie waarborgt dat de student de procedurele stappen niet alleen uitvoert, maar dit ook doet met de nodige precisie die nodig is in een klinische omgeving \cite{Pottle2019, Alaker2016VR}.

\section{Implementatie van het dataloggingsysteem}
Een essentieel onderdeel van de ontwikkeling is het opstellen van objectieve evaluatiegegevens, een behoefte die vaak wordt onderstreept in de literatuur over medische XR-training \cite{Radianti2020VRHE, Kyaw2019VR}. Er is een gespecialiseerde Logging Manager ingesteld die zich onafhankelijk van de visuele representatie richt. Elke interactie, beslissing en tijdelijke variabele wordt vastgelegd met een nauwkeurigheid die tot op de milliseconde reikt. Dit systeem hanteert conceptueel de richtlijnen van de Experience API , zodat leerervaringen op een gestructureerde manier kunnen worden opgeslagen en uitgewisseld. \cite{xAPIADL}

Tijdens de sessie worden de verzamelde gegevens opgeslagen om de CPU-belasting te verlagen en na het einde van het scenario worden ze geserialiseerd naar een JSON-formaat. Deze aanpak stelt ons in staat om belangrijke prestatie-indicatoren, zoals de tijd tot de eerste schok of de pauzes tussen de acties, op een objectieve manier te onderzoeken. Deze datagedreven benadering levert een aanzienlijke voordelen op in vergelijking met conventionele observatietechnieken, omdat het mogelijk maakt om patronen in fouten te herkennen en de effectiviteit van de simulatie als leermiddel te bewijzen \cite{Hung2024VR, Padilha2019VR}.

\section{Optimalisatie en technische beperkingen}
Tijdens het ontwikkelingsproces zijn de unieke beperkingen van standalone XR-hardware in overweging genomen. Om een hoge mate van 'presence' te waarborgen zonder de cognitieve belasting negatief te beïnvloeden, is de visuele omgeving aangepast aan de principes van multimedia learning \cite{Mayer2009Multimedia, WitmerSinger1998}. Het verminderen van ingewikkelde schaduwen en het toepassen van geoptimaliseerde meshes garandeert dat de grafische processor genoeg capaciteit heeft voor de verwerking van de fysica en handtracking \cite{Slater2009PlaceIllusion}.

Hoewel de huidige technologie nog beperkingen heeft met betrekking tot tactiele feedback, is er geprobeerd dit technisch te verhelpen met behulp van geavanceerde visuele en auditieve feedbackmechanismen. De compressies worden beïnvloed door een auditieve metronoom, terwijl de interactie met de AED wordt ondersteund door nauwkeurige gesproken instructies die de werking van het werkelijke apparaat nabootsen.\cite{ILCOR2023Immersive} Deze combinatie van technische en educatieve aspecten creëert een leeromgeving die, ondanks de hardwarematige beperkingen, een hoge mate van procedurele nauwkeurigheid garandeert (Fowler 2015VR, Shorey Ng2021VRSR).